% answers_short.tex —— 2007 年·全国大纲卷Ⅱ·理科 · 简短答案速查
% 由 tools/make_short_answers.py 自 answers.tex 的 \gksolhead/\gkanswer 提取，
% 并与 2007/2007_数学真题_全国2理_答案要点.md 逐题比对后生成，勿手改（改源 MD 后重跑）。

\begin{gknotice}
\textbf{使用说明}：本节只给最终答案便于快速核对，完整推导见解析版。答案为编写者独立推导，\textbf{非官方标准答案}。
\end{gknotice}

\gkpart{选择题}
\begin{center}\begin{tabular}{ccccccccccccc}
\toprule
题号 & 1 & 2 & 3 & 4 & 5 & 6 & 7 & 8 & 9 & 10 & 11 & 12 \\
\midrule
答案 & B & C & C & D & A & C & A & A & C & B & B & B \\
\bottomrule
\end{tabular}\end{center}

\gkpart{填空题}
\begin{itemize}[leftmargin=2.6em]
  \item[\textbf{13.}] $-42$
  \item[\textbf{14.}] $0.8$
  \item[\textbf{15.}] $2+4\sqrt2$
  \item[\textbf{16.}] $-\dfrac52$
\end{itemize}

\gkpart{解答题}
\begin{itemize}[leftmargin=2.6em]
  \item[\textbf{17.}] (1) $y=2\sqrt3+6\sin x+2\sqrt3\cos x$，$0<x<\dfrac{2\pi}3$；(2) 最大值为 $6\sqrt3$
  \item[\textbf{18.}] (1) $p=\dfrac15$；(2) $P(\xi=0)=\dfrac{316}{495}$，$P(\xi=1)=\dfrac{32}{99}$，$P(\xi=2)=\dfrac{19}{495}$
  \item[\textbf{19.}] (1) 见解析版；(2) 二面角 $A-EF-D$ 满足 $\cos\theta=\dfrac5{\sqrt{105}}$
  \item[\textbf{20.}] (1) 圆为 $x^2+y^2=4$；(2) $\overrightarrow{PA}\cdot\overrightarrow{PB}\in\left(-\dfrac45,\dfrac45\right)$
  \item[\textbf{21.}] (1) $a_n=1+(a_1-1)\left(-\dfrac12\right)^{n-1}$；(2) $b_n<b_{n+1}$
  \item[\textbf{22.}] (1) 切线为 $y=(3t^2-1)x-2t^3$；(2) 见解析版
\end{itemize}
